\section{Methodology}
\label{method}
\noindent Building upon the multi-granular annotations in SDA, we propose the Multi-granularity Relation Alignment (MRA) framework. 
As depicted in Fig.~\ref{fig: framework}, MRA is optimized through two complementary constraints:
1) fine-grained alignment between image regions and description phrases via Region-Phrase Contrastive (RPC) Learning and Region-Phrase Matching (RPM) Learning; 
2) coarse-grained alignment between full images and texts through Image-Text Alignment.




\subsection{Region-Phrase Contrastive Learning}
\noindent Region-Phrase Contrastive (RPC) Learning aligns the region-phrase pair $(I_R, T_R)$ with contrastive learning.
Fig.~\ref{fig: framework} shows that the $N$ image-text pairs and the $N$ corresponding region-phrase pairs are fed into MRA in a training batch $D$.
For $N$ pairs of region and phrase $\{(I_R^1, T_R^1), (I_R^2, T_R^2), ...\}^N$, $(I_R^j, T_R^j)$ is a positive match, while $(I_R^j, T_R^{\bar{j}})$ is a negative match ($j, \bar{j} \in [1, N], j \neq \bar{j}$). 
We first obtain the region embeddings $E_V(I_R)$ and phrase embeddings $E_T(T_R)$ through the Vision Encoder $E_V$ and the Text Encoder $E_T$.


\noindent\textbf{Region Encoding.}
Following the existing work~\cite{yang2023towards}, the Swin Transformer (Swin-Base)~\cite{liu2021swin} is deployed as the Vision Encoder $E_V$. Feeding the region to $E_V$ is a straightforward strategy for region encoding; however, this fashion is not friendly to the small-sized region. In addition, the sole region lacks the surrounding contexts for better representation. With this consideration, we resort to indexing the region features from the full image features. 
Formally, given the region $I_R$, its corresponding image $I$ is evenly divided into $N^I$ non-overlapping patches (with a patch size of $32 \times 32$) of fixed size first. 
The image size is set to $224 \times 224$, and $N^I = 49$.
Then, the sequence of $N^I$ patches is fed into the Transformer block layers of $E_V$, outputting $N^I$ high-dimensional embeddings.
Region embedding $E_V(I_R)$ derives from $V$: $E_V(I_R)= \{v^{i} | \text{ if } v^{i} \text{ overlaps with } I_R, i=1,2,\dots,N^I\}$.
Then the mean feature of $E_V(I_R)$ is calculated and used as the \texttt{[CLS]} patch embedding of $I_R$, which is together with the patch features in $E_V(I_R)$ to form the full representation of the region $I_R$.


\noindent\textbf{Phrase Encoding.}
Without loss of generality, we use the first 6 layers of BERT as the Text encoder $E_T$ to extract the textual representation. 
Specifically, for a given phrase $T_R$, it is first tokenized by a BERT Tokenizer with a vocabulary size of $30,522$. 
A \texttt{[CLS]} token is added at the beginning of the set of tokens.
The tokens are then passed into the six Transformer layers of BERT to obtain the textual embedding $E_T(T_R)$.


The obtained region and phrase embeddings are first averaged separately and mapped to the low-dimensional region features by different MLPs. 
Then, we predict the similarity between regions and phrases by calculating the cosine similarity $s(\cdot,\cdot)$ between region and phrase. 
Given a pair of region and phrase $(I_R, T_R)$, the region-to-phrase similarity is defined as follows:
\begin{equation}
\label{eq: S_R2P}
S_{\text{R2P}} = \frac{\exp(s(I_R, T_R)/\tau)}{\sum_{i=1}^{N}\exp(s(I_R, T_R^i)/\tau)},
\end{equation}
where we still use $I_R, T_R$ to represent the respective features for brevity, and $\tau$ is a learnable temperature parameter.
Similarly, the phrase-to-region similarity is $S_{\text{P2R}}$.
Finally, the RPC loss function is defined as:
\begin{equation}
\label{eq: L_PRC}
\begin{split}
\mathcal{L}_{\text{RPC}} = -\frac{1}{2} \mathbb{E}[\log S_{\text{R2P}} + \log S_{\text{P2R}}].
\end{split}
\end{equation}


\subsection{Region-Phrase Matching Learning.}
\noindent Region-Phrase Matching (RPM) Learning is a binary classification task used to further differentiate between positive (matched) and negative (unmatched) region-phrase pairs. 
If a negative sample is randomly selected to form a negative sample pair, the classification problem would be too easy for the model. 
Therefore, based on RPC, we form negative sample pairs by selecting one hard-negative sample for each region and phrase through a Hard-Negative Sampling strategy. 
Specifically, for a given region $I_R$, a Hard-Negative Sample $\Bar{T}_R$ is randomly selected in proportion to the similarity score in Eq. \ref{eq: S_R2P}. 
Similarly, a hard-negative sample for each phrase can be obtained analogously. 
Then, the $N$ matched region-phrase pairs and $2N$ corresponding mismatched region-phrase pairs are fed into the Fusion Encoder $E_F$ together to obtain the fusion embeddings, as shown in Fig. \ref{fig: framework}.

\noindent\textbf{Fusion Encoder $E_F$.}
We adopt the last six layers of BERT as $E_F$. 
$E_F$ takes phrase embedding as inputs and fuses region embedding at each cross-attention module of the Transformer Blocks. 
Specifically, the phrase embedding is used as a query (q in Fig. \ref{fig: framework}), and the region embedding works as key (k) and value (v). 
The fusion embedding yielded by the Fusion Encoder is used to perform further cross-modal alignment tasks.
We project the \texttt{[CLS]} embeddings of the fused embeddings into the two-dimensional space by MLP to get the predicted matching likelihood $\hat{p}_m(I_R, T_R)$.

The Region-Phrase Matching (RPM) loss is obtained by calculating the cross-entropy loss between the predicted matching likelihood and the ground-truth matching likelihood $p_m(I_R, T_R)$:
\begin{equation}
\label{eq: L_PRM}
\begin{split}
\mathcal{L}_{\text{RPM}} =& - \mathbb{E}[p_{m}(I_R, T_R)\log \hat{p}_{m}(I_R, T_R)\\
&+ (1-p_{m}(I_R, T_R))\log (1-\hat{p}_{m}(I_R, T_R))],
\end{split}
\end{equation}
\noindent where $p_{m}$ is a ground-truth binary label. If the image and text are true-matches, $p_{m} = 1$, otherwise $p_{m} = 0$.


\noindent\textbf{Discussion. Why do we need region-phrase level alignment?}
While recent methods for image-text matching have productively incorporated both global and local feature representations \cite{jiang2023cross, bai2023rasa, yang2023towards, yan2022clip}, a systematic, explicit alignment between discrete image regions and their corresponding textual phrases remains largely unexplored. This constitutes a significant limitation, as fine-grained cross-modal understanding fundamentally requires the precise grounding of visual constituents within their descriptive linguistic contexts. Our work directly bridges this gap by introducing an explicit region-phrase alignment mechanism, facilitated by the novel structured annotations in our proposed SDA dataset.
As shown in Table~\ref{tab: abl}, Method M1, which relies solely on global image-text alignment, exhibits a marked performance shortfall relative to our complete MRA framework. This comparative analysis confirms that integrating region-phrase alignment contributes a substantial and quantifiable improvement to retrieval accuracy.
Furthermore, our investigation into even finer-grained object-word alignment (variants 1 and 2 in Table~\ref{tab: abl}) yields a critical observation: exceeding an optimal level of granularity can be counterproductive. The suboptimal performance of object-word alignment suggests that excessive fragmentation of the representation may introduce noise and disrupt the compositional semantics preserved by phrase-level descriptions. This evidence highlights that region-phrase alignment achieves an effective balance, providing the necessary discriminative detail to surpass global methods while avoiding the loss of contextual coherence inherent in over-segmented approaches.




\begin{figure*}[!t]
\centering
\begin{minipage}[t]{0.45\linewidth}
    \centering
    \small    \makeatletter\def\@captype{table}\makeatother
    \caption{Performance comparisons with state-of-the-art methods on the CUHK-PEDES dataset.}
    \label{tab:sota_CUHK}
    \centering
    \resizebox{1\linewidth}{!}{ 
        \renewcommand\arraystretch{1} 
        \begin{tabular}[h!]{p{2.6cm}|m{1.2cm}<{\centering}m{1.2cm}<{\centering}m{1.2cm}<{\centering}m{1.2cm}<{\centering}}
            \shline
            \textbf{Method} & \textbf{R@1} & \textbf{R@5} & \textbf{R@10} & \textbf{mAP}  \\
            \hline
            CNN-RNN~\cite{reed2016learning} & 8.07 & - & 32.47 & - \\ 
            GNA-RNN~\cite{li2017person} & 19.05 & - & 53.64 & - \\ 
            PWM-ATH~\cite{chen2018} & 27.14 & 49.45 & 61.02 & - \\ 
            GLA~\cite{chen2018improving} & 43.58 & 66.93 & 76.2 & - \\
            Dual Path~\cite{zheng2020dual} & 44.40 & 66.26 & 75.07 & - \\
            CMPM+CMPC~\cite{zhang2018deep} & 49.37 & - & 79.21 & - \\
            MIA~\cite{niu2020improving} & 53.10 & 75.00 & 82.90 & -\\
            A-GANet~\cite{liu2019deep} & 53.14 & 74.03 & 81.95 & - \\  
            ViTAA~\cite{wang2020vitaa} & 55.97 & 75.84 & 83.52 & 51.60 \\ 
            IMG-Net~\cite{wang2020img} & 56.48 & 76.89 & 85.01 & - \\
            CMAAM~\cite{aggarwal2020text} & 56.68 & 77.18 &	84.86 & - \\
            HGAN~\cite{zheng2020hierarchical} &	59.00 & 79.49 & 86.62 & - \\
            DSSL~\cite{zhu2021dssl} & 59.98 & 80.41 & 87.56 & - \\
            MGEL~\cite{wang2021text} & 60.27 & 80.01 & 86.74 & - \\
            SSAN~\cite{ding2021semantically} & 61.37 & 80.15 & 86.73 & - \\
            TBPS~\cite{han2021text} & 61.65 &80.98 & 86.78 & - \\
            TIPCB~\cite{chen2022tipcb} & 63.63 &82.82 & 89.01 & - \\
            LBUL~\cite{wang2022look} & 64.04 &82.66 & 87.22 & - \\
            CAIBC~\cite{wang2022caibc} & 64.43 & 82.87 &88.37 & - \\
            AXM-Net~\cite{farooq2022axm} & 64.44 & 80.52 & 86.77 & 58.73 \\
            SRCF~\cite{suo2022simple} & 64.88 &83.02 & 88.56 & - \\
            LGUR~\cite{shao2022learning} & 65.25 & 83.12 & 89.00 & - \\
            CFine~\cite{yan2022clip} & 69.57 & 85.93 & 91.15 & - \\
            IRRA~\cite{jiang2023cross} & 73.38 & 89.93 & 93.71 & 66.13 \\
            SAMC~\cite{lu2024mind} & 74.03 & 89.18 & 93.31 & 68.42 \\
            RaSa~\cite{bai2023rasa} & 76.51 & 90.29 & 94.25 & \textbf{69.38} \\
            APTM~\cite{yang2023towards} & 76.53	& 90.04 & 94.15 & 66.91\\
            \hline
            Baseline & 71.23 & 87.54 & 92.43 & 64.75 \\
            Ours & \textbf{77.21}	& \textbf{90.66} & \textbf{94.46} & 68.50 \\
            \shline
        \end{tabular}
    }    
\end{minipage}
\hspace{0.4cm}
\quad
\begin{minipage}[t]{0.45\linewidth}
    \centering
    \begin{minipage}[t]{1.0\linewidth}
        \centering
        \small
        \makeatletter\def\@captype{table}\makeatother
        \caption{Performance comparisons with state-of-the-art methods on the ICFG-PEDES dataset.} 
        \label{tab:sota_ICFG}
        \centering
        \resizebox{1\linewidth}{!}{
            \renewcommand\arraystretch{1} 
            \begin{tabular}[h!]{p{3cm}|m{1.2cm}<{\centering}m{1.2cm}<{\centering}m{1.2cm}<{\centering}m{1.2cm}<{\centering}} 
                \shline
                \textbf{Method} & \textbf{R@1} & \textbf{R@5} & \textbf{R@10} & \textbf{mAP}  \\
                \hline
                Dual Path~\cite{zheng2020dual}   & 38.99 & 59.44 & 68.41 & - \\
                CMPM+CMPC~\cite{zhang2018deep}   & 43.51 & 65.44 & 74.26 & - \\
                MIA~\cite{niu2020improving}      & 46.49 & 67.14 & 75.18 & - \\
                SCAN~\cite{lee2018stacked}       & 50.05 & 69.65 & 77.21 & - \\
                ViTAA~\cite{wang2020vitaa}       & 50.98 & 68.79 & 75.78 & - \\
                SSAN~\cite{ding2021semantically} & 54.23 & 72.63 & 79.53 & - \\
                IVT~\cite{shu2022see}            & 56.04 & 73.60 & 80.22 & - \\
                LGUR~\cite{shao2022learning}     & 59.02 & 75.32 & 81.56 & - \\
                CFine~\cite{yan2022clip}         & 60.83 & 76.55 & 82.42 & - \\
                IRRA~\cite{jiang2023cross}       & 63.46 & 80.25 & 85.82 & 38.06 \\
                SAMC~\cite{lu2024mind} & 63.68 & 79.69 & 85.21 & \textbf{42.41} \\
                RaSa~\cite{bai2023rasa}          & 65.28 & 80.40 & 85.12 & 41.29 \\
                APTM~\cite{yang2023towards}      & 68.51 & 82.99 & 87.56 & 41.22\\
                \hline
                Baseline & 66.02 & 81.67 & 86.78 & 39.04 \\ 
                Ours & \textbf{68.93} & \textbf{83.46} & \textbf{88.07} & 41.38 \\
                \shline
            \end{tabular}
        }
    \end{minipage}    
    \begin{minipage}[t]{1.0\linewidth}
        \centering
        \small
        \makeatletter\def\@captype{table}\makeatother
        \caption{Performance comparisons with state-of-the-art methods on the RSTPReid dataset.}
        \label{tab:sota_RSTP} 
        \centering
        \resizebox{1\linewidth}{!}{
            \renewcommand\arraystretch{1} 
            \begin{tabular}[h!]{p{3cm}|m{1.2cm}<{\centering}m{1.2cm}<{\centering}m{1.2cm}<{\centering}m{1.2cm}<{\centering}}
                \shline
                \textbf{Method} & \textbf{R@1} & \textbf{R@5} & \textbf{R@10} & \textbf{mAP}  \\
                \hline
                DSSL~\cite{zhu2021dssl}     & 32.43 & 55.08 & 63.19 & - \\
                LBUL~\cite{wang2022look}    & 45.55 & 68.20 & 77.85 & - \\
                IVT~\cite{shu2022see}       & 46.70 & 70.00 & 78.80 & - \\
                CAIBC~\cite{wang2022caibc}  & 47.35 & 69.55 & 79.00 & - \\
                CFine~\cite{yan2022clip}    & 50.55 & 72.50 & 81.60 & - \\
                IRRA~\cite{jiang2023cross}  & 60.20 & 81.30 & 88.20 & 47.17 \\
                SAMC~\cite{lu2024mind} & 60.80 & 82.35 & 89.00 & 49.67 \\
                RaSa~\cite{bai2023rasa}     & 66.90	& \textbf{86.50} & 91.35 & 52.31 \\
                APTM~\cite{yang2023towards} & 67.50	& 85.70 & \textbf{91.45} & 52.56 \\
                \hline
                Baseline & 56.30 & 77.55 & 85.20 & 45.88 \\
                Ours & \textbf{68.15} & 86.30 & 91.10 & \textbf{53.77}\\
                \shline
            \end{tabular}
        }
    \end{minipage}  
\end{minipage}  
\end{figure*}


\subsection{Image-Text Alignment}
\noindent In addition to region-phrase level alignment (RPC and RPM), we combine the tasks of Image-Text Contrastive (ITC) learning, Image-Text Matching (ITM) learning, and Masked Language Modeling (MLM) to constrain the model to align pedestrian images and text descriptions.


\noindent\textbf{Image-Text Contrastive (ITC) Learning} is performed in a similar fashion of RPC. The ITC loss $\mathcal{L}_{\text{ITC}}$ is obtained following Eq.\ref{eq: S_R2P}-\ref{eq: L_PRC}.


\noindent\textbf{Image-Text Matching (ITM) Learning.}
The goal of Image-Text Matching (ITM) Learning is to predict whether the image-text pairs are positive pairs or not.
The ITM loss $\mathcal{L}_{\text{ITM}}$ is computed similarly using Eq. \ref{eq: L_PRM}.







\noindent\textbf{Masked Language Modeling (MLM).}
Masked Language Modeling (MLM) predicts masked words using information from matched image-text pairs as clues. 
In this paper, we randomly mask out text tokens with a probability of $25\%$. 
Among the masked tokens: 
1) $10\%$ are replaced by random text tokens; 
2) $80\%$ are replaced by special text tokens \texttt{[MASK]}; 
3) the remaining $10\%$ stay unchanged.
Given a matched $(I, T)$, mask $T$ as $\hat{T}$ and use $\hat{f}^i$ to denotes fused embedding of masked token $\hat{t}^i$.
We predict what the token $\hat{t}^i$ should be by a classified MLP header, denoted as $\hat{p}_{\text{mask}}^i(I, \hat{T}) = \text{Softmax}(\text{MLP}(\hat{f}^i))$.
$p_{\text{mask}}^i$ is a one-hot vector representing the ground-truth probability. Then the MLM loss can be formulated as follows by the cross-entropy:
\begin{equation}
\label{eq: L_MLM}
\mathcal{L}_{\text{MLM}} = -\mathbb{E}[p_{\text{mask}}^i(I, \hat{T}) log( \hat{p}_{\text{mask}}^i(I, \hat{T}))].
\end{equation}
MRA jointly optimizes Region-Phrase Alignment (RPC and RPM) and Image-Text Alignment. 
The overall MRA loss is:
\begin{equation}
\label{eq: L_MRA}
\mathcal{L}_{\text{MRA}} = \mathcal{L}_{\text{ITC}} + \mathcal{L}_{\text{ITM}} +\mathcal{L}_{\text{MLM}} + \beta (\mathcal{L}_{\text{RPC}} + \mathcal{L}_{\text{RPM}}),
\end{equation}
\noindent where $\beta$ denotes the weight of RPC and RPM loss.


